%--------------------------------------------
% PREAMBLE (to be loaded before each paper)
%--------------------------------------------


% -------------------------------------PACKAGES
%\usepackage{graphicx}      % include this line if your document contains figures
%\usepackage[square, comma, numbers]{natbib}        % required for bibliography
%\bibliographystyle{alphadin}
%\usepackage{bibtex}


%user added
%	a) basic
\usepackage{amsmath, amssymb, amsfonts, mathtools}%amsthm, ... symbols & Co. %%TO INSTALL mathtools, bbding, csquotes, ulem, xparse
%\usepackage{pifont, bbding}  %symbols
%	\let\theoremstyle\relax
\usepackage{csquotes}
\usepackage[normalem]{ulem} %underlying, without replacing em
\usepackage[ngerman,english]{babel}
%\usepackage[]{mcode} %matlab
\usepackage{enumerate}
\usepackage{xparse} %define fancy commands (http://stackoverflow.com/questions/1812214/latex-optional-arguments)


%\usepackage[utf8]{inputenc}
%\usepackage[T1]{fontenc}

%	b) nice to haves

\usepackage{color}
\definecolor{TUMblau}		{RGB}{0, 101, 189}
\definecolor{TUMgruen}		{RGB}{162, 173, 0}
\definecolor{TUMorange}		{RGB}{227, 114, 34}
\definecolor{TUMdunkelrot}	{RGB}{156,013,022}
\definecolor{gray}			{RGB}{140,140,140}


\usepackage[hyperfootnotes=false,
pdffitwindow=false,						% do not resize document window to fit document size
pdftitle={},							% title
pdfauthor={Maria Cruz Varona},			% author
%pdfcreator={LaTeX},					% creator of the document
linktocpage=true,						% true: only pages in the table of contents is made into a link
colorlinks=true,						% false: boxed links; true: colored links
linkcolor=TUMblau,						% color of internal links
citecolor=TUMdunkelrot,					% color of links to bibliography
urlcolor=TUMgruen,						% color of external links
draft = true
]{hyperref}

\hypersetup{
	pdfinfo={
		Keywords={}
	}
}


\usepackage{todonotes} %[colorinlistoftodos, disable]
\usepackage{booktabs,units} %pmat
%\usepackage[justification=centering]{caption}
%\usepackage{url}

\usepackage{algorithm, algpseudocode} %algpseudocode
	\renewcommand{\algorithmicrequire}{\textbf{Input:}}
	\renewcommand{\algorithmicensure}{\textbf{Output:}}
	

%	c) TIKZ
\usepackage{pgfplots}
	% and optionally (as of Pgfplots 1.3):
	\pgfplotsset{compat=newest} %newest
%	\pgfplotsset{compat=1.10} %previously 1.9
	\pgfplotsset{plot coordinates/math parser=false}
	\pgfplotsset{yticklabel style={text width=2.5em,align=right}}
	\pgfplotsset{/pgf/number format/.cd, fixed,precision=3}
	
	%	TIKZ-EXTERNALIZE
%	\usepgfplotslibrary{external}
%		\tikzexternalize[prefix=./pics/externaltikz/] %shell escape=-enable-write18] %optimize=false,
		
	%	---- README TIKZ -----
	%
	%	In order to get tikzexternalize to work, make sure that
	%	- your main file name has no empty spaces
	%	- the conflict with todonotes is resolved!
	%	----------------------
			
	%	Remove compatibility issues with todonotes!
%	\newcommand{\todos}[1]{\tikzexternaldisable\todo[inline]{#1}\tikzexternalenable}
		%	\makeatletter		
		%	\makeatother
			\newcommand{\todos}[1]{\todo[inline]{#1}}
		
	% use the following lines to produce pdf and png with ImageMagick (install Ghostscript and ImageMagick for this to work)
%	\tikzset{
%		png export/.style={
%			external/system call={
%				pdflatex \tikzexternalcheckshellescape -halt-on-error -interaction=batchmode -jobname "\image" "\texsource"&&
%				convert -density 300 -transparent white "\image.pdf" "\image.png"
%			}
%		}
%	}
%	\tikzset{png export}
		
	%	Remake all pictures
%			\tikzset{external/force remake}
	
	\newlength\myheight
	\newlength\mywidth
%		\setlength\mywidth{85 mm}
%		\setlength\myheight{35 mm} 
	

	%	\usetikzlibrary{plotmarks}
	%	Define standardized lengths for tikz plots throughout the thesis
%	\newlength{\subplotsfheight}
%	\newlength{\subplotsfwidth}
%	\newlength{\plotsfheight}
%	\newlength{\plotsfwidth}
%	\newlength{\appsubplotsfheight}
%	\newlength{\appsubplotsfwidth}
%	\setlength{\appsubplotsfheight}{4cm}
%	\setlength{\appsubplotsfwidth}{6.5cm}
%	\setlength{\subplotsfheight}{3cm}
%	\setlength{\subplotsfwidth}{5cm} 
%	\setlength{\plotsfheight}{5cm}
%	\setlength{\plotsfwidth}{6.5cm} 
	




% *********************** USER DEFINED ENVIRONMENTS (amsthm package)**************************************
\newtheorem{df}{\bf Definition}
%\newtheorem{lm}{Lemma}
%\newtheorem{remark}{Remark}
%\newtheorem{thm1}{Theorem}
%\newtheorem{crl}{Corollary}
%\newtheorem{propos}{Proposition}


% % USER DEFINED FUNCTIONS
% 1) Symbols 
\newcommand{\Reals}{\mathbb{R}}
\newcommand{\Complex}{\mathbb{C}}
\newcommand{\Htwo}{\mathcal{H}_2}
\newcommand{\Hinf}{\mathcal{H}_\infty}
\renewcommand{\th}{^{\text{th}}}
\newcommand{\BigO}[1]{\mathcal{O}\left(#1\right)}

\newcommand{\defeq}{\vcentcolon=}
\newcommand{\logicand}{\wedge}
\newcommand{\logicor}{\vee}

% 2) Matrix operations
\newcommand{\trans}[1]{#1^\top}
\newcommand{\htrans}[1]{#1^H}
\newcommand{\hermit}{^H}
\newcommand{\inv}[1]{#1^{-1}}
\newcommand{\invt}[1]{#1^{-\top}}
\newcommand{\pd}{\succ}
\newcommand{\nd}{\prec}


% 3) Operators
\newcommand\abs[1]{\left|#1\right|}
\DeclareMathOperator{\vecspan}{span}
	\newcommand\vspan[1]{\vecspan\left(#1\right)}
	\newcommand\vecspace[1]{\mathcal{#1}}
\DeclareMathOperator{\diagmat}{diag}
	\newcommand\diag[1]{\diagmat\left(#1\right)}
\DeclareMathOperator{\matrixrank}{rank}
	\newcommand\rank[1]{\matrixrank #1}
\DeclareMathOperator{\matrixrang}{rang}
	\newcommand\rang[1]{\matrixrang\left(#1\right)}

\newcommand{\norm}[1]{\left\lVert#1\right\rVert}
\newcommand{\diff}[2]{\frac{\mathrm{d}\, #1}{\mathrm{d}\, #2}}
\newcommand{\innProd}[2]{\left\langle#1, \; #2\right\rangle}

% 4) Formatting
\newcommand{\ts}{\!} %negative space for formulas in text
\newcommand{\nix}{\mbox{ }} %empty space



%---- MOR comamnds
\newcommand{\krylov}[3]{\vecspace{K}_{#1}\left(#2, #3\right)}
\newcommand{\shift}{\sigma}
\newcommand{\As}{A_{\shift}}
\newcommand{\Asx}{A-\shift E}
\newcommand{\inputkrylov}{\krylov{n}{\inv{\As} E}{\inv{\As}B}}
\newcommand{\outputkrylov}{\krylov{n}{\invt{\As} \trans{E}}{\invt{\As} \trans{C}}}
\newcommand{\moment}{C\left(\inv{\As}E\right)^i \inv{\As} B}

\newcommand{\pencil}{\mu E - A}

\newcommand{\GramC}{\mathcal{W}_c}
\newcommand{\GramO}{\mathcal{W}_o}

%----- theorems
%\newtheorem{thm}{Theorem}[section]
%\newtheorem{lemma}{Lemma}
%%\newtheorem{corollary}{Corollary}
%\newtheorem{proposition}{Proposition}

%\theoremstyle{definition}
%\newtheorem{definition}{Definition}

%\theoremstyle{remark}
%\newtheorem{remark}{Remark}